\subsection{Resultados para QFT}

% A diferencia de Grover, QFT mostró un comportamiento fuertemente dependiente de la estructura
% del estado de entrada. Esta subsección es particularmente importante porque permite delimitar
% el alcance real del esquema propuesto: cuando la entrada conserva regularidades, la
% compresión puede ser útil; cuando predomina una estructura de alta entropía, esa utilidad
% desaparece e incluso puede invertirse. Por ello, QFT constituye el caso más ilustrativo para
% discutir las condiciones bajo las cuales la compresión del vector de estado resulta
% conveniente.

\subsubsection{Superposición periódica}

\begin{table}[H]
    \centering
    \small
    \setlength{\tabcolsep}{3pt}
    \caption{QFT con superposición periódica}
    \label{tab:qft-pattern-results}
    \resizebox{\linewidth}{!}{%
    \begin{tabular}{lcccccc}
        \hline
        \textbf{Qubits} & \textbf{Estrategia} & \textbf{Mem. pico (MB)} & \textbf{Ahorro mem. (\%)} & \textbf{Tiempo (s)} & \textbf{Tiempo rel. (x)} & \textbf{Error} \\
        \hline
        22 & Referencia & 73.6 & 0.0 & 0.44 & 1.00 & 0 \\
        22 & Blosc & 35.7 & 51.4 & 1.52 & 3.48 & 0 \\
        22 & ZFP & 19.0 & 74.2 & 2.52 & 5.75 & $3.79\times 10^{-11}$ \\
        23 & Referencia & 139.6 & 0.0 & 1.21 & 1.00 & 0 \\
        23 & Blosc & 34.5 & 75.3 & 3.17 & 2.61 & 0 \\
        23 & ZFP & 25.4 & 81.8 & 5.68 & 4.68 & $3.79\times 10^{-11}$ \\
        24 & Referencia & 271.6 & 0.0 & 2.75 & 1.00 & 0 \\
        24 & Blosc & 92.4 & 66.0 & 6.99 & 2.54 & 0 \\
        24 & ZFP & 45.6 & 83.2 & 12.10 & 4.41 & $3.79\times 10^{-11}$ \\
        \hline
    \end{tabular}
    }
\end{table}

% Cuando la entrada presenta periodicidad, ambas estrategias comprimidas logran reducciones de
% memoria apreciables. Blosc ahorra entre 51.4\% y 75.3\% y mantiene un costo temporal entre
% $2.54\times$ y $3.48\times$ respecto a la referencia. ZFP incrementa aún más el ahorro, hasta
% el rango 74.2\%--83.2\%, pero lo hace a costa de una penalización temporal mayor, entre
% $4.41\times$ y $5.75\times$, e introduce un error máximo absoluto constante del orden de
% $10^{-11}$.

% Estos resultados muestran que, en QFT, la compresión puede ser útil si el estado de entrada
% conserva una organización suficientemente regular. Sin embargo, también evidencian que la
% elección entre Blosc y ZFP depende del criterio de optimización dominante. Si se privilegia la
% exactitud exacta junto con un compromiso más equilibrado entre memoria y tiempo, Blosc resulta
% más conveniente. Si se prioriza el mayor ahorro espacial y se acepta una pequeña discrepancia
% numérica, ZFP puede ofrecer una ventaja puntual.

\subsubsection{Superposición de alta entropía}

\begin{table}[H]
    \centering
    \small
    \setlength{\tabcolsep}{3pt}
    \caption{QFT con superposición de alta entropía}
    \label{tab:qft-random-results}
    \resizebox{\linewidth}{!}{%
    \begin{tabular}{lcccccc}
        \hline
        \textbf{Qubits} & \textbf{Estrategia} & \textbf{Mem. pico (MB)} & \textbf{Ahorro mem. (\%)} & \textbf{Tiempo (s)} & \textbf{Tiempo rel. (x)} & \textbf{Error} \\
        \hline
        22 & Referencia & 71.6 & 0.0 & 0.77 & 1.00 & 0 \\
        22 & Blosc & 85.3 & -19.0 & 9.67 & 12.52 & 0 \\
        22 & ZFP & 73.1 & -2.1 & 79.52 & 103.01 & $3.53\times 10^{-11}$ \\
        23 & Referencia & 135.7 & 0.0 & 2.00 & 1.00 & 0 \\
        23 & Blosc & 152.2 & -12.1 & 20.98 & 10.48 & 0 \\
        23 & ZFP & 187.6 & -38.2 & 171.47 & 85.65 & $2.82\times 10^{-11}$ \\
        24 & Referencia & 263.6 & 0.0 & 4.33 & 1.00 & 0 \\
        24 & Blosc & 392.8 & -49.0 & 44.30 & 10.24 & 0 \\
        24 & ZFP & 409.8 & -55.5 & 376.40 & 87.03 & $3.79\times 10^{-11}$ \\
        \hline
    \end{tabular}
    }
\end{table}

% La situación cambia por completo cuando la entrada presenta alta entropía. En este caso, ni
% Blosc ni ZFP ofrecen una ventaja real frente al almacenamiento no comprimido. El ahorro de
% memoria se vuelve negativo en todos los tamaños, lo que significa que ambas estrategias
% consumen más memoria que la referencia. En Blosc, ese sobrecosto espacial oscila entre 12.1\%
% y 49.0\%; en ZFP, alcanza hasta 55.5\%.

% El deterioro temporal también es pronunciado. Blosc incrementa el tiempo entre
% $10.24\times$ y $12.52\times$, mientras que ZFP lo lleva a factores entre $85.65\times$ y
% $103.01\times$. Aunque los errores de ZFP permanecen en un rango pequeño, del orden de
% $10^{-11}$, el problema central ya no es la precisión admitida, sino la falta de
% compresibilidad del estado.

% Desde el punto de vista experimental, este es el resultado más importante de QFT. La
% compresión del vector completo no puede asumirse como una estrategia generalmente beneficiosa:
% cuando el estado carece de regularidad local o repetición suficiente, la compresión deja de
% reducir memoria y además introduce sobrecostos temporales elevados. En consecuencia, la
% utilidad del esquema depende no solo del circuito en abstracto, sino de la estructura
% concreta del estado que ese circuito debe representar.
